\begin{textsample}{POS Dim 3 – go – Score 20.00 – go/go\_inf\_22.txt}  \label{ex:f3_pos_049}
3.4.2.
Organização da Ação Pedagógica É importante demarcar que o professor é um agente de cultura, responsável por \textbf{apoiar} a \textbf{criança} na constituição da sua identidade e na compreensão do mundo físico e sociocultural.
Quanto a compreensão, uso e apropriação da linguagem oral e escrita o professor tem a função de mediar e ampliar as experiências da \textbf{criança} relacionadas a esse processo.
Para tanto, orienta -se que observe quais são os interesses e \textbf{curiosidades} das \textbf{crianças}, o que as mobiliza, quais os assuntos de suas conversas, quais \textbf{brincadeiras} prefere, se há algum conflito na turma, o que provoca essa situação e como resolvê -la.
No caso de \textbf{bebês}, que ainda não conseguem verbalizar, é preciso estar atento às suas outras linguagens, sorrisos, choros, \textbf{balbucios} para identificar seus interesses e necessidades.
No trabalho com a linguagem oral destaca -se que a \textbf{criança} aprende a falar e a pensar nas interações.
Portanto, precisa ser incentivada a falar com os outros, dar recados, \textbf{contar} fatos, histórias, expressar sentimentos, \textbf{desejos}, emoções, opiniões.
O ato de comunicar \textbf{ajuda} a \textbf{criança} a pensar sobre tudo que vivencia, desenvolvendo sua fala e seu pensamento.
Assim, é necessário o professor buscar estratégias de \textbf{escuta} dessa \textbf{criança} e planejar situações que oportunizem a construção da linguagem por meio de diálogos, de contação e leitura de histórias, de \textbf{brincadeiras}, de exploração de objetos e materiais, de interações com outras \textbf{crianças} e \textbf{adultos}.
Para auxiliar a \textbf{criança} no processo de nomeação de objetos, sentimentos e ações é fundamental organizar o espaço com objetos que permitam \textbf{brincar} de esconder/aparecer, afastar/buscar, subir/descer, \textbf{empilhar}, montar/desmontar, oportunizando -lhe experimentar essas diferentes ações.
Também, é fundamental a experiência de \textbf{brincar} com outras \textbf{crianças}, promovendo proximidade física com parceiros de idades próximas, ou não.
Ações como ofertar objetos, partilhar o olhar, se \textbf{interessar} pelo outro são comportamentos sociais primordiais que possibilitam a emergência do diálogo, da negociação, do conflito, da troca de \textbf{afetos}, da fala.
Cabe ao professor identificar quais \textbf{brincadeiras}, jogos, canções, parlendas e versos a \textbf{criança} já sabe para que possa introduzir novos.
Jogos linguísticos como trava-línguas, adivinhas tradicionais, lengalengas, provérbios, canções, rimas e poemas ampliam a participação da \textbf{criança} na cultura oral, possibilitando reflexões que contribuirão para a compreensão da escrita.
Ressalta -se a importância de planejar e desenvolver atividades envolvendo conversas com a \textbf{criança}, realizadas em diferentes contextos em que a linguagem oral se faz presente, tais como a roda de conversa, as \textbf{brincadeiras} livres, especialmente as de faz de conta, os debates na sala sobre assuntos relevantes e de interesses da \textbf{criança}.
Assim, é possível ampliar suas possibilidades de falar sobre a vida, seus \textbf{desejos}, seus medos, seus sonhos, enfim, sobre o que veem e pensam sobre o mundo.
Para tanto, faz -se importante usar gêneros orais variados, tais como : piadas, notícias, recontos, causos, relatos, informes, recados etc.
Outras possibilidades são as leituras em voz alta de fábulas, textos instrucionais, contos, histórias, adivinhas, biografias, bilhetes, dentre várias outras opções de textos escritos.
O objetivo é oferecer acesso a diferentes formas de estruturação de textos escritos e aos seus conteúdos.
Diferente da leitura em voz alta, as narrativas orais também são fundamentais para a ampliação e a compreensão do escrito e da expressão linguística, pois possibilitam a análise e a criação de conceitos, quando há mediação do professor.
A \textbf{criança}, ao ser convidada a \textbf{contar} uma história a partir de ilustrações, é mobilizada a organizar a estrutura narrativa com o apoio das imagens.
Uma outra estratégia é incentivá -la a lembrar de uma história por meio de perguntas, isso também auxilia no processo de organizá -la numa sequência lógica e compreensível ao ouvinte, bem como a pensar nas palavras e expressões utilizadas pelo autor da história.
A narrativa oral, por meio de recontos, auxilia a \textbf{criança} a organizar mentalmente o resumo de um texto, a parafrasear ou repetir a história utilizando ora palavras do texto escutado, ora suas próprias palavras, fazer comentários sobre o seu conteúdo ou até mesmo reformular a história.
A roda de conversa é outro dispositivo de diálogo e de construção de narrativas a ser desenvolvida regularmente e se constitui num momento privilegiado de interação entre \textbf{adultos} e \textbf{crianças}, \textbf{crianças} e \textbf{crianças}, favorecendo a aprendizagem da \textbf{escuta} do outro, da fala, da argumentação, da expressão de ideias e da construção de significados coletivos.
Para que a conversa seja ato de linguagem e de diálogo é indispensável a alternância e o respeito às falas, na perspectiva de garantir o envolvimento e a participação de todos na pronúncia, \textbf{escuta}, contraposição, aceitação ao posicionamento do outro, perguntas e respostas.
As práticas de leitura e as formas de ler, dependem das suas intenções e finalidades, do tipo de livro, das conversas e interações que antecedem, que acompanham e que sucedem a leitura.
Para realizá -la com a \textbf{criança} ou grupo de \textbf{crianças}, é importante o professor se atentar ao movimento de repetição muito próprio delas como abrir e fechar o livro várias vezes, rever ilustrações e pedir para repetir trechos da história ou ainda dizer : “ \textbf{Conta} de novo? ”.
Também são comuns os \textbf{gestos} imitativos das \textbf{crianças} ao reproduzir expressões corporais, entonações de vozes, ritmo de leitura presenciadas por elas.
A forma como a \textbf{criança} lê está estreitamente relacionada com a performance do professor ao dar corpo e voz ao texto.
A performance está relacionada à entonação que o professor usa ao ler, ao destaque que é dado ao pronunciar certas palavras, aos \textbf{gestos} que acompanham a leitura, todos esses elementos são fundamentais para que a \textbf{criança} consiga produzir sentido ao que é lido.
Por isso o professor tem que ler para e com a \textbf{criança}, ler de corpo inteiro, vivenciar o texto com ela, se deixar levar pela emoção, pela imaginação.
Mas, o que ler para a \textbf{criança}?
O que motiva as escolhas do livro literário?
Como o ambiente de leitura é organizado?
Quais formas de ler são possíveis no espaço da instituição?
Essas reflexões não podem vir desvinculadas de uma concepção de \textbf{criança} capaz, ativa, curiosa e interessada em conhecer o mundo, a vida.
Muito menos desarticulada do Projeto Político-Pedagógico da instituição.
O professor, ao escolher o que ler para a \textbf{criança}, deve assegurar, entre outros aspectos, que a \textbf{criança} tenha acesso a livros de qualidade literária e estética, que utilizam materiais, técnicas e recursos gráficos variados, tais como : desenhos, fotografias, massinhas, gravuras, colagens etc.
A presença dos livros na instituição não é suficiente para garantir que a \textbf{criança} tenha acesso a eles.
É preciso organizar o espaço de forma a incentivar e promover as leituras, com locais fixos, como salas ou cantinhos de leitura, armários para guardar o acervo literário, prateleiras, entre outros, ou formas alternativas que possibilitem os livros transitarem entre salas e as \textbf{crianças}, como, cestas, baús, caixas, \textbf{carrinhos} de mão.
Quanto a escrita, o espaço da instituição educacional deve se constituir num ambiente que instigue a \textbf{curiosidade} da \textbf{criança}, despertando e aguçando nela o \textbf{desejo} e a vontade de ler e escrever, a partir de situações comunicativas reais, vivenciadas no cotidiano.
Para tanto, serão elencadas algumas ações pedagógicas, dentre várias possíveis, que favorecem a participação da \textbf{criança} nas culturas escritas.
São elas : - colocar placas indicativas nas portas - cozinha, refeitório, sanitário feminino, masculino ou acessível - e ensinar a \textbf{criança} a identificar ; - organizar murais informativos para as famílias e/ou responsáveis com diferentes gêneros textuais e atualizá -los com regularidade ; - escrever e ler junto com a turma os bilhetes que são enviados para casa ; - produzir textos coletivos da turma, sendo o professor o escriba ; - permitir que as \textbf{crianças}, desde \textbf{bebês}, tenham acesso a materiais que possibilitem deixar marcas -giz de cera, \textbf{lápis} de cor, carvão, tinta guache etc. ; - incentivar as \textbf{crianças} maiores, a escreverem de forma espontânea - rabiscos, garatujas, letras ; - promover sarau de poesias, festival de contadores de histórias, etc.
Conclui -se afirmando que a \textbf{criança} da Educação \textbf{Infantil} pode ( e deve ) ser convidada a ler e a escrever, mesmo sem ter o domínio da leitura no sentido restrito.
Ela é capaz de interpretar sinais gráficos, relacionar imagens e textos, antecipar e inferir sentidos, atentar -se a organização de textos em diferentes suportes ou contextos de escrita, a refletir sobre o funcionamento da língua, observando as possibilidades de combinações entre letras e \textbf{sons}, sentidos e significados.

% matched lemmas: adulto, afeto, ajuda, apoiar, balbucio, bebê, brincadeira, brincar, carro, contar, criança, curiosidade, desejo, empilhar, escuta, gesto, infantil, interessar, lápis, som
\end{textsample}
